\section{Simplicity of the Alternating Groups}

Recall that a conjugacy class in \(S_n\) consists of the set of all elements of a fixed cycle type.

\begin{example}
    For example, all the \(\para{2, 3}\)-cycles in \(S_5\) form a conjugacy class.
\end{example}

\begin{theorem}[Simplicity of \(A_n\) for \(n \geq 5\)]
    \label{thm:an-simple}%
    Let \(n \geq 5\). Then \(A_n\) is simple.
\end{theorem}

\begin{proof}
    There are three steps of the proof:
    \begin{enumerate}
        \item \(A_n\) is generated by \(3\)-cycles;
        \item Any \(H \normalsub A_n\) that contains a \(3\)-cycle contains all of them;
        \item Any non-trivial \(H \normalsub A_n\) does contain a \(3\)-cycle.
    \end{enumerate}

    It should be clear that the three statements combined gives the desired statement. Note in particular that the second statement is not trivial -- since we are working in \(A_n\), not \(S_n\).

    \begin{proof}[of 1]
        Any \(g \in A_n\) is a product of evenly-many transpositions. We consider pairs of transpositions, then this gives three possibilities:
        \begin{itemize}
            \item \(\perm{ab} \perm{ab} = e\).
            \item \(\perm{ab} \perm{bc} = \perm{abc}\) is a \(3\)-cycle.
            \item \(\perm{ab} \perm{cd} = \perm{acb} \perm{acd}\) is a product of \(3\)-cycles.
        \end{itemize}
    \end{proof}

    \begin{proof}[of 2]
        Any two \(3\)-cycles are conjugate when viewed in \(S_n\). Let \(\delta, \delta'\) be \(3\)-cycles, then we may write
        \[
            \delta' = \sigma \delta \sigma\inverse
        \]
        for \(\sigma \in S_n\). But observe that \(n \geq 5\), so there is some transposition \(\tau \in S_n\) disjoint from \(\delta\), which means it commutes with \(\delta\). Therefore,
        \[
            \delta' = \para{\sigma \tau} \delta \para{\tau\inverse \sigma\inverse} = \para{\sigma \tau} \delta \para{\sigma \tau}\inverse
        \]
        and hence \(\sigma \tau \in A_n\), showing \(\delta\) and \(\delta'\)are conjugate in \(A_n\).
    \end{proof}
\end{proof}